%% ===========================================================================
%% Controlled-substrate results for the SCAFF audit, restructured around the five-step
%% empirical narrative: (1) observe natural divergence D_s^nat, (2) attribute it as direct
%% (D_s^A) or inherited (D_s^Z), (3) apply the prohibited-dependence premise (a normative
%% judgment call, stated but not measured), (4) measure harm only where a reference Omega_s
%% exists, (5) repair the responsible stage and test whether the downstream disparity
%% disappears. Elicit/Retrieve/Rank are the pipeline's empirical core under this substrate;
%% Explain and Memory are demoted to their own "Extensions" subsection below, since a
%% single-turn endpoint audit is blind to them by construction, not merely underpowered.
%%
%% This file is now the sole surviving results section for BOTH local lineages. It was
%% originally written against the notation of the Overleaf project's main_yd.tex (Yashar's
%% SCAFF rewrite): D_s^{nat} / D_s^A / D_s^Z for the three coordinates, \stage{} for stage
%% names, and the equation labels eq:natural, eq:crossovercell, eq:direct, eq:state,
%% eq:violation, eq:repair. main_yd.tex is not present in this local checkout (Overleaf-only),
%% so paper/section_scaff.tex -- the local theory section, written with the equivalent macros
%% \Dnat/\DA/\DZ and labels eq:crossover/eq:dnat/eq:ddirect/eq:dinherited -- now carries alias
%% labels (eq:crossovercell/eq:natural/eq:direct/eq:state) on its own equations so the
%% cross-references below resolve locally. eq:violation, eq:quality and eq:repair have no
%% standalone equation in section_scaff.tex (they name a decision rule, a harm quantity and a
%% procedure that are stated in prose there), so references to them below point to
%% Section~\ref{sec:scaff} in prose rather than to a nonexistent equation label.
%% paper/section_scaff_results.tex -- the counterpart this file used to duplicate for the
%% fair_trace_acm.tex/main.tex lineage -- has been retired per the note that used to live in
%% it, now that this file is \input directly by fair_trace_acm.tex.
%%
%% Provenance of every number: the two-domain controlled prototype
%% (FAIR_TRACE_SCAFF_TwoDataset_Toy_Colab.ipynb; the crossover numbers were independently
%% reproduced in 0817_FAIR-TRACE_SCAFF_crossover_summary.ipynb), seed 42, 24 scenarios x 4
%% repeats per domain, LLM judge disabled. Re-executed locally; the re-execution reproduced the
%% notebook's recorded outputs exactly. Figure regenerated by scripts/make_fig_scaff.py (the
%% `paper' variant, drawn 1:1 at ACM text width; v1--v4 are the same content at presentation
%% size).
%%
%% Section~\ref{sec:controlled-robustness} is NOT from the prototype: it is an independent
%% re-implementation (scripts/scaff_robustness.py, driven by scripts/run_scaff_robustness.py,
%% narrated in 0819_FAIR-TRACE_SCAFF_localization_robustness.ipynb). Every number it quotes is
%% emitted to paper/scaff_robustness_results.json by that driver; re-running is byte-identical.
%% Do not merge its numbers into Table~\ref{tab:controlled-coords} -- different substrate.
%%
%% \input this file from Section~\ref{sec:results} of fair_trace_acm.tex (or, on Overleaf, from
%% main_yd.tex immediately before \bibliographystyle, if/when that lineage is reconciled).
%% ===========================================================================

\providecommand{\NDCG}{\mathrm{NDCG}}

\section{Results on a Controlled Substrate}
\label{sec:controlled}

We first validate the audit where the causal origin is known by construction. Two controlled
domains---movie and restaurant recommendation---are built so that the latent user profile, the
set of acceptable actions, item relevance and memory support are all exactly known, and so that
a protected-attribute dependence is \emph{planted} at a single designated stage. Each domain
contributes $24$ scenarios at $4$ repeats, distributed uniformly over the planted stage
(including control scenarios with no plant); a stage is exposed to the descriptor only when it is
that scenario's planted stage. Every stage output passes the validity gate of
Section~\ref{sec:controlled-validity} below before any fairness quantity is computed on it. The
agent is implemented as transparent stage functions rather than through an agent framework,
because the crossover of Equation~\eqref{eq:crossovercell} requires each parent state $Z_s$ to be
addressable, holdable and replayable, and requires it to be provable that a stage did not observe
the descriptor---three properties a framework that owns state passing and sampling internally
does not expose (Section~\ref{sec:scaff}).

\paragraph{Scope of the planted dependence}
The planted faults are unit tests of the measurement instrument. They are not estimates of
real-world demographic disparity and assert nothing about the preferences or behaviour of any
group. What they establish is whether the audit recovers a dependence whose location we already
know---a precondition for trusting it where we do not.

\begin{table}[t]
  \caption{Mean per-stage coordinates on the controlled substrate ($24$ scenarios
    $\times\,4$ repeats per domain), computed on one primary coordinate per stage.
    \emph{noise} is the within-condition repeat distance and bounds the effect size that
    $D_s^{A}$ must exceed to be interpretable at all.}
  \label{tab:controlled-coords}
  \begin{tabular}{llrrrr}
    \toprule
    Domain & Stage & $D_s^{\mathrm{nat}}$ & $D_s^{A}$ & $D_s^{Z}$ & noise \\
    \midrule
    Movies      & \stage{Elicit}   & 0.056 & 0.056 & 0.000 & 0.000 \\
                & \stage{Retrieve} & 0.205 & 0.136 & 0.069 & 0.166 \\
                & \stage{Rank}     & 0.396 & 0.094 & 0.302 & 0.328 \\
                & \stage{Explain}  & 0.225 & 0.067 & 0.158 & 0.008 \\
                & \stage{Memory}   & 0.156 & 0.056 & 0.101 & 0.000 \\
    \midrule
    Restaurants & \stage{Elicit}   & 0.056 & 0.056 & 0.000 & 0.000 \\
                & \stage{Retrieve} & 0.220 & 0.146 & 0.074 & 0.120 \\
                & \stage{Rank}     & 0.490 & 0.156 & 0.333 & 0.370 \\
                & \stage{Explain}  & 0.197 & 0.072 & 0.125 & 0.014 \\
                & \stage{Memory}   & 0.160 & 0.056 & 0.104 & 0.004 \\
    \bottomrule
  \end{tabular}
\end{table}

\subsection{Step 1: Observing Natural Divergence}
\label{sec:controlled-step1}

Table~\ref{tab:controlled-coords} reports $D_s^{\mathrm{nat}}$, the diagonal contrast of
Equation~\eqref{eq:natural}, for every stage before any attribution is attempted: this is what a
stage-wise comparison of the two counterfactual trajectories sees on its own, and it is where any
audit that only replays the pipeline twice and diffs the outputs necessarily stops. The pattern is
immediate at \stage{Rank}, which carries the largest natural divergence of any stage in both
domains ($D_s^{\mathrm{nat}} = 0.396$ movies, $0.490$ restaurants), well above \stage{Retrieve}
($0.205$, $0.220$), \stage{Explain} ($0.225$, $0.197$) and \stage{Memory} ($0.156$, $0.160$). Read
alone, Equation~\eqref{eq:natural} cannot distinguish a stage that responded to the descriptor
from one that faithfully inherited an already-corrupted upstream state; it only reports that the
two counterfactual arms differ by the time execution reaches that stage. An audit that stopped
here would direct remediation effort at whichever stage happens to sit furthest downstream of the
true fault, which---absent further information---is close to the opposite of the right basis for
assigning responsibility. Step 2 resolves this by decomposing $D_s^{\mathrm{nat}}$ into the direct
and inherited coordinates of Equations~\eqref{eq:direct} and~\eqref{eq:state}.

\subsection{Step 2: Attributing Direct vs.\ Inherited Dependence}
\label{sec:controlled-step2}

\subsubsection{Localization}
\label{sec:rq1}

Selecting, per scenario, the stage with the largest $D_s^{A}$---and \emph{no stage} when none
exceeds the decision threshold---recovers the planted stage in all $48$ scenarios across the two
domains, including the controls in which nothing was planted. \stage{Elicit}-, \stage{Retrieve}-
and \stage{Rank}-only injections are each localized without confusion, and the confusion matrix
is diagonal.

A bare $100\%$ is not, however, the reportable form of this result, and
Section~\ref{sec:controlled-robustness} re-derives it: on an independent re-implementation at
$120$ scenarios the figure carries a $95\%$ Wilson interval of $[0.969, 1.000]$, and across five
seeds it runs from $0.958$ to $1.000$. The exact $100\%$ was in part seed luck. The margin used
here was also never calibrated, and calibrating it against the control scenarios' own noise
changes the answer materially; that is taken up in the same section.

One qualification survives unaddressed. Accuracy is measured at a single planted effect size, and
the plants here are strong. The regime that matters for a deployed audit is the weak-effect
regime, so the fully honest form of this result is accuracy as a function of planted effect size
relative to the noise floor, which the controlled substrate can supply but has not yet been run to
produce.

Table~\ref{tab:controlled-coords} also shows where the instrument is weakest: at \stage{Rank}
the within-condition repeat distance ($0.328$, $0.370$) exceeds the direct effect being estimated
($0.094$, $0.156$). The primary coordinate for that stage is a top-one gap, a binary quantity that
flips under decoding stochasticity while remaining blind to a genuine reordering that leaves the
head of the list intact---simultaneously noisy and insensitive, which is the worst combination for
a detection statistic. Section~\ref{sec:controlled-robustness} resolves this, and the resolution
has two parts rather than one. A repeat-noise floor of that size next to a direct effect of that
size is the signature of a crossover whose exogenous noise was \emph{not} held fixed across the
four cells, contrary to Equation~\eqref{eq:crossovercell}; holding it fixed removes the
contamination entirely. Where the agent's sampling cannot be controlled, the coordinate then
decides whether the audit works at all: the top-one gap separates planted from unplanted
\stage{Rank} faults at an AUC of $0.566$, against $0.804$ for a difference in $\NDCG@5$. The
\stage{Rank} row of Table~\ref{tab:controlled-coords} should be read as superseded by
Table~\ref{tab:rank-coords}.

\subsubsection{Natural Divergence Overattributes}
\label{sec:rq2}

The overattribution flagged in Step 1 is visible in Table~\ref{tab:controlled-coords} at
\stage{Rank}, which carries the largest natural divergence of any stage in both domains
($D_s^{\mathrm{nat}} = 0.396$ and $0.490$) while most of that divergence is inherited
($D_s^{Z} = 0.302$, $0.333$) rather than direct ($D_s^{A} = 0.094$, $0.156$). An audit reading
Equation~\eqref{eq:natural} alone would direct remediation at the ranker.
\stage{Retrieve} shows the opposite composition---direct exceeding inherited in both
domains---and is where responsibility is correctly assigned.

Figure~\ref{fig:scaff} in Section~\ref{sec:scaff} shows the mechanism on a single scenario whose
fault is planted at \stage{Elicit}. On the $A=a$ arm, \stage{Elicit} adds the preference
\texttt{tone=romantic}, which the latent profile does not support. Natural divergence is then
non-zero at four consecutive downstream stages---\stage{Retrieve} $0.20$, \stage{Rank} $0.75$,
\stage{Explain} $0.50$, \stage{Memory} $0.33$---and reading Equation~\eqref{eq:natural} per stage
implicates all of them. Holding each stage's parent state fixed and swapping only the descriptor
gives $D_s^{A} = 0$ at every one of the four, with the entire divergence carried by $D_s^{Z}$.
Only \stage{Elicit} has $D_s^{A} > 0$. The overattribution is therefore not a matter of degree:
four stages that never observed the descriptor are reported as attribute-sensitive by natural
divergence and exonerated by Equation~\eqref{eq:direct}.

One reporting consequence follows from Table~\ref{tab:controlled-coords} and is easy to get
wrong. \stage{Elicit}'s corpus mean $D_s^{A}$ is $0.056$, below the decision margin, even though
\stage{Elicit} is the planted stage in a sixth of scenarios and is localized in every one of
them. A corpus mean over per-scenario coordinates summarizes a distribution and is not itself a
verdict; the diagnosis rule of Section~\ref{sec:scaff} (``Reading the coordinates'') is defined
per scenario, and the corpus-level object to report is the distribution of per-scenario
diagnoses.

\subsubsection{Process Sensitivity Masked at the Endpoint}
\label{sec:rq3}

Table~\ref{tab:controlled-masking} cross-tabulates whether a pair is process-sensitive---some
stage with $D_s^{A}$ above the margin---against whether its final recommendation differs beyond
the endpoint-equivalence margin. Of $151$ process-sensitive pairs, $77$ are endpoint-equivalent:
for just over half of the trajectories in which some stage responded directly to the descriptor,
the delivered recommendation is indistinguishable. Per domain, $36.5\%$ (movies) and $43.8\%$
(restaurants) of all pairs are masked in this sense. These are precisely the cases an
output-level consumer-fairness audit records as fair, and they are the empirical argument for
attributing the process (Step 2) rather than stopping at the endpoint.

\begin{table}[t]
  \caption{Process sensitivity against endpoint difference, pooled over both domains. The $77$
    pairs in the upper right are process-sensitive trajectories that an output-only audit scores
    as fair.}
  \label{tab:controlled-masking}
  \begin{tabular}{lrr}
    \toprule
     & endpoint differs & endpoint equivalent \\
    \midrule
    process sensitive             & 74 & 77 \\
    no direct process sensitivity &  0 & 41 \\
    \bottomrule
  \end{tabular}
\end{table}

Part of this masking is structural rather than incidental, and the structural half---the
contribution of \stage{Explain} and \stage{Memory}---is separated out and reported in
Section~\ref{sec:controlled-extensions} rather than pooled in here, because it is a claim of a
different and stronger kind than the modest power gap the remaining stages show.

The empty lower-left cell should be read as a property of this substrate rather than a general
guarantee: a planted fault is the only route by which the two arms can differ here, so no
endpoint difference can arise without some direct sensitivity. Under a stochastic agent that cell
will be populated, and it should be reported rather than assumed empty.

\subsection{Step 3: The Prohibited-Dependence Premise}
\label{sec:controlled-step3}

Steps 1 and 2 are measurements: $D_s^{\mathrm{nat}}$, $D_s^{A}$ and $D_s^{Z}$ are computed from
trace data and require no judgment about what a stage is entitled to depend on. Whether a
non-zero $D_s^{A}$ constitutes \emph{procedural unfairness} is a separate, normative claim,
supplied by the auditor rather than measured (Section~\ref{sec:scaff}, ``Three claims,
separated''), and it is stated per stage because the same descriptor can be a legitimate input to
one stage and a prohibited one to another---a safety filter and a preference-inference stage do
not carry the same obligation with respect to it.

On this controlled substrate the premise is fixed by construction rather than argued: the planted
stage is, in every scenario, stipulated to be the illegitimate consumer of the descriptor, so a
non-zero $D_s^{A}$ at the planted stage is by design a procedural violation, and a non-zero
$D_s^{A}$ elsewhere would not be, were it to occur. This is a property of the unit-test design, not
a result about which stages are legitimately sensitive in a deployed agent. We did not run a
version of the substrate in which the premise itself is varied---for instance, treating
\stage{Elicit}'s use of a user-volunteered preference correlated with $A$ as legitimate while
treating the same dependence at \stage{Retrieve} or \stage{Rank} as prohibited---and we do not
report one here; deciding and validating stage-specific prohibition premises on a realistic agent
is future work, not a result this substrate establishes.

\subsection{Step 4: Measuring Harm Where a Reference Exists}
\label{sec:controlled-step4}

The strongest of the three claims separated in Section~\ref{sec:scaff} is disparate harm: that the
affected party is worse off, which requires a reference $\Omega_s$ against which one output can be
judged better than another. On this substrate a reference is available by construction---the
latent user profile is exact---so harm can be reported wherever direct sensitivity is found and
left unreported otherwise. On the scenario of Figure~\ref{fig:scaff} the divergence is not merely
a difference: the latent profile gives \texttt{tone = cerebral}, so the preference
\stage{Elicit} added on the $A=a$ arm is wrong against that reference, and the quality difference
$\Delta Q_s$ (the harm quantity of Section~\ref{sec:scaff}) is identified rather than merely
observed.

Where a scenario provides no defensible reference for a stage, harm is recorded as \emph{not
identified} and the sensitivity result stands on its own; we take that as the correct output
rather than a gap to be filled in with an LLM judge (Section~\ref{sec:scaff}). Step 5 below
reports several such not-identified cells for the repair experiment specifically---removability is
a harm-adjacent quantity in the same sense, and the same discipline governs it: report the gap
rather than manufacture a number to fill it.

\subsection{Step 5: Repair and Retest}
\label{sec:controlled-step5}

\subsubsection{Repair Removes Downstream Divergence}
\label{sec:rq4}

Applying the repair procedure of Section~\ref{sec:scaff}---replacing the implicated stage's output
with its valid counterpart and re-running all descendants---reduces final ranking divergence by
$0.576$ (movies) and $0.588$ (restaurants). The same operation at a stage carrying inherited
divergence only reduces it by $0.170$ and $0.187$. The attribution therefore predicts the effect
of acting on it, which a coordinate that merely described divergence would not, and the gap
between the two numbers is the substantive content of this test.

One constraint on the experiment has to be stated because violating it manufactures a spurious
success. Removability must not be measured on the output of the stage being repaired: a repair at
\stage{Rank} scores a removability exactly equal to its own pre-repair divergence in both domains
($0.249$ and $0.278$), because substituting the ranked list substitutes the very quantity the
outcome is computed on. That figure measures the definition of the metric, not the effect of the
repair. Removability is therefore evaluated strictly downstream of the repaired stage, which for a
repair at \stage{Rank} in a single-turn setting means on the following turn.
Table~\ref{tab:repair-matrix} shows how far the violation propagates, and it is further than the
\stage{Rank} cell alone: scored on the ranking, a repair at \stage{Retrieve} also removes $0.939$
of an \stage{Elicit}-planted divergence, because \stage{Retrieve}'s output largely determines the
ranking. The numbers quoted in the previous paragraph inherit that confound and are superseded by
the matrix, presented with the robustness checks in Section~\ref{sec:controlled-robustness} below.

\subsection{Extensions: \stage{Explain} and \stage{Memory}}
\label{sec:controlled-extensions}

The five-step narrative above is run to completion only for \stage{Elicit}, \stage{Retrieve} and
\stage{Rank}, which we treat as the paper's empirical core. \stage{Explain} and \stage{Memory} are
demoted to extensions rather than dropped, and the reason is structural rather than a matter of
statistical power: both stages act \emph{downstream} of the ranked list, so within a single turn
no fault planted at either one can change the endpoint under any parameterisation. A single-turn
endpoint audit is therefore \emph{blind by construction} to \stage{Explain} and \stage{Memory}
faults, not merely underpowered against them, and no amount of additional repetition at this
substrate's effect size would fix that: the limitation is where the observable sits in the
pipeline, not how noisy the measurement of it is.

Section~\ref{sec:controlled-robustness} verifies the structural claim directly rather than arguing
it: splitting the masking pairs of Section~\ref{sec:controlled-step2} by whether the endpoint
\emph{could} move, the $40$ pairs planted at \stage{Explain} or \stage{Memory} are all
process-sensitive and all endpoint-equivalent, with a maximum endpoint distance of \emph{exactly}
$0.0$---compared with a real but modest $5.0\%$ masking rate among the $80$ pairs planted at
\stage{Elicit}, \stage{Retrieve} or \stage{Rank}. The repair experiment of Step 5 shows the same
asymmetry from a different angle: Table~\ref{tab:repair-matrix} does not even carry a column for
repairing \stage{Memory}, because no stage lies downstream of it within a turn against which
removability could be scored at all, and removability for \stage{Rank}, \stage{Explain} and
\stage{Memory} faults is reported as \emph{not identified} in a single turn rather than filled in
(Step 4's discipline applied to a second quantity).

The blindness is a property of the single-turn horizon, not of the stages themselves, and we have
run---though not fully integrated into the five-step narrative above---one experiment that
demonstrates this. Under an explicit carry-over model in which the agent's stored facts and the
reason tags it gave are echoed into the next turn's revealed preferences, every planted stage moves
the turn-two endpoint, including \stage{Explain} ($0.408$) and \stage{Memory} ($0.163$), while the
controls remain at exactly $0.000$. This result is reported in
Section~\ref{sec:controlled-robustness} below because it comes from the same robustness
re-implementation, not from a full second pass of Steps 1--5 at \stage{Explain} and
\stage{Memory}; running that pass on a genuinely multi-turn substrate, rather than the single
follow-up turn probed here, is future work and is what would be required to promote these two
stages back into the paper's empirical core.

\subsection{Validity and Testability}
\label{sec:controlled-validity}

The validity precondition rejected $2.1\%$ of \stage{Explain} outputs in both domains---all of
them a mismatch between the item named in the explanation and the top-ranked item, and all
concentrated on one scenario per domain. Those runs are excluded from \stage{Explain}'s
coordinates rather than scored, and the rejection rate is reported beside them: a stage that
frequently emits structurally invalid output is a finding in its own right, and averaging over
such runs silently would report it as invariance.

\subsection{Robustness of the Instrument}
\label{sec:controlled-robustness}

The results above are the prototype's own run. Four properties of that run are load-bearing and
none of them was established: the choice of \stage{Rank}'s coordinate, the decision margin, the
precision of the localization figure, and the scoring of the repair experiment. This subsection
establishes them on an independent re-implementation of the measurement---one domain rather than
two, transparent stage functions, \texttt{numpy} and \texttt{pandas} only---parameterised to sit
in the same regime as the prototype rather than to reproduce it exactly ($120$ scenarios at $8$
repeats, catalog of $60$ items, five seeds; under the original top-one coordinate its
\stage{Retrieve} repeat-noise is $0.181$ against the prototype's $0.166$, its \stage{Rank} direct
effect $0.104$ against $0.094$, its \stage{Rank} natural divergence $0.376$ against $0.396$).
Re-execution is byte-identical, and every number below is emitted to a machine-readable results
file rather than transcribed.

Figure~\ref{fig:scaff-robustness} summarises the three findings that change how the results above
should be read.

\paragraph{The exogenous noise must be shared across crossover cells}
Equation~\eqref{eq:crossovercell} writes the cell as
$Y_s^{i,j} = \phi_s(f_s(Z_s^{i}, A{=}j; \epsilon))$ with a single $\epsilon$ common to all four
cells. This is not a notational convenience. $D_s^{A}$ is a difference between two cells, so if
each cell draws its own randomness that difference is a mixture of the effect and the sampling
noise and $D_s^{A}$ inherits the stage's entire noise floor; when $\epsilon$ is held fixed the
noise cancels in the difference and $D_s^{A}$ is exact. The prototype's reported \stage{Rank} row,
a noise coordinate of $0.328$ beside a direct effect of $0.094$, is what the second regime
produces. Running the identical trajectories both ways, \stage{Rank}'s three coordinates move
from $(D^{\mathrm{nat}}, D^{A}, D^{Z}) = (0.376, 0.104, 0.272)$ under a shared $\epsilon$ to
$(0.850, 0.762, 0.788)$ under independent draws---indistinguishable from each other and from the
noise floor, so the decomposition carries no information at that stage at all. Localization falls
from $100\%$ to between $16.7\%$ and $82.5\%$ depending on the coordinate.

This is the highest-leverage correction in this section and it costs nothing to apply. It also
imposes a requirement on the evaluation settings of Section~\ref{sec:controlled-limits}: either
the agent's sampling is seed-controllable, or the audit must run on a coordinate that survives
resampling.

\paragraph{\stage{Rank}'s coordinate}
Table~\ref{tab:rank-coords} scores five candidate distances on the same trajectories in the
noisy regime, where the choice actually matters. \emph{effect} is the mean $D_s^{A}$ at
\stage{Rank} on \stage{Rank}-planted scenarios; \emph{floor} is the residual $D_s^{A}$ where
nothing was planted, which is the quantity a threshold must clear. The top-one gap detects a
planted fault at an AUC of $0.566$, which is close enough to chance to be unusable, and being
binary it is simultaneously blind to a true reordering that leaves the head of the list intact.
A difference in $\NDCG@5$ against the latent profile reaches $0.804$.

\begin{table}[t]
  \caption{Candidate coordinates for \stage{Rank}, scored on identical trajectories in the
    independent-$\epsilon$ regime. \emph{floor} is the residual $D_s^{A}$ on scenarios with
    nothing planted; \emph{AUC} is the separation of \stage{Rank}-planted scenarios from all
    others by $D_s^{A}$; \emph{loc.} is end-to-end localization accuracy over all five stages.
    Under a shared $\epsilon$ the floor is $0$ and all five localize at $100\%$, so the ranking
    below is the price of not controlling the agent's sampling.}
  \label{tab:rank-coords}
  \begin{tabular}{lrrrr}
    \toprule
    coordinate & effect & floor & AUC & loc. \\
    \midrule
    $\Delta\NDCG@5$ (needs $\Omega_s$) & 0.043 & 0.030 & \textbf{0.804} & \textbf{82.5\%} \\
    top-3 Jaccard                      & 0.565 & 0.517 & 0.684 & 34.2\% \\
    Kendall $\tau$                     & 0.526 & 0.483 & 0.680 & 33.3\% \\
    RBO ($p = 0.9$)                    & 0.452 & 0.421 & 0.659 & 32.5\% \\
    top-one gap (prototype)            & 0.791 & 0.756 & 0.566 & 16.7\% \\
    \bottomrule
  \end{tabular}
\end{table}

The best coordinate is the one that requires an oracle, which is a real cost rather than a
recommendation: $\Delta\NDCG@K$ is available only where $\Omega_s$ is defensible. The operating
rule adopted here is $\Delta\NDCG@K$ where a reference exists and RBO where none does, and the
top-one gap in neither case.

\paragraph{The decision margin}
The margin $\tau = 0.10$ was set by hand. It can be derived instead: the control scenarios have
nothing planted, so every $D_s^{A}$ they exhibit is instrument noise, and an upper quantile of
their per-scenario maximum $D_s^{A}$ fixes $\tau$ at a stated false-positive budget. With
$\epsilon$ shared the controls show \emph{exactly} zero direct sensitivity, so the calibrated
margin is $\tau^\star = 0$ and any strictly positive $D_s^{A}$ is signal.

The hand-set value is not merely arbitrary; it silently encodes an assumption about the
coordinate's scale, and that assumption fails for the coordinate this section recommends.
$\Delta\NDCG@5$ effects live around $0.04$, entirely beneath $0.10$, so at the hand-set margin
every \stage{Rank} plant is missed and accuracy caps at $83.3\%$---on all five seeds, without
exception. Adopting a graded coordinate while retaining the inherited margin would therefore have
made the audit worse, not better.

\paragraph{Localization, with an interval}
Under a shared $\epsilon$, the calibrated margin and $\Delta\NDCG@5$, localization is correct on
$120$ of $120$ scenarios. The bootstrap interval is $[1.000, 1.000]$, which is the familiar
degeneracy of resampling at a boundary proportion and should not be reported as evidence; the
Wilson score interval, $[0.969, 1.000]$, is the honest one. Across five seeds accuracy runs from
$0.958$ to $1.000$ with a mean of $0.992$, so the exact $100\%$ was in part seed luck. The
defensible statement of this result is \emph{$0.958$--$1.000$ across seeds, $95\%$ Wilson interval
$[0.969, 1.000]$ at $n = 120$}, and the calibrated margin beats the hand-set one on every seed.

\paragraph{Repair, scored without circularity}
Table~\ref{tab:repair-matrix} makes the constraint stated in Section~\ref{sec:controlled-step5}
measurable. Scored the original way---every cell on the final ranking distance---the entire
\stage{Rank} column reads $1.0$: repairing the ranker reports complete removal of the disparity
for a fault planted at \stage{Elicit}, at \stage{Retrieve} or at \stage{Rank} alike, so the
experiment cannot distinguish them. Restricting the outcome to \stage{Rank} alone does not repair
this, because \stage{Retrieve}'s output largely determines the ranking: a \stage{Retrieve} repair
still removes $0.939$ of an \stage{Elicit}-planted divergence despite \stage{Retrieve} being
faithful. Only scoring strictly downstream of the repaired stage separates them, and it does so
cleanly---diagonal $1.0$, largest off-diagonal entry $0.20$.

\begin{table}[t]
  \caption{Fraction of divergence removed by repairing each stage (columns) for a fault planted
    at each stage (rows). Left: scored on the final ranking distance, as originally. Right:
    scored strictly downstream of the repaired stage. A trustworthy experiment puts mass on the
    diagonal only. \emph{n.i.} marks cells where removability is \emph{not identified}: the
    repaired stage has no downstream divergence to remove. No column exists for repairing
    \stage{Memory}: within a turn, no stage lies downstream of it against which removability
    could be scored at all.}
  \label{tab:repair-matrix}
  \begin{tabular}{lrrr@{\hskip 2em}rrrr}
    \toprule
    & \multicolumn{3}{c}{scored on the ranking}
    & \multicolumn{4}{c}{scored strictly downstream} \\
    \cmidrule(r){2-4}\cmidrule(l){5-8}
    planted at & \stage{Eli.} & \stage{Ret.} & \stage{Rank}
               & \stage{Eli.} & \stage{Ret.} & \stage{Rank} & \stage{Exp.} \\
    \midrule
    \stage{Elicit}   & \textbf{1.00} & 0.94 & 1.00 & \textbf{1.00} & 0.20 & 0.00 & 0.00 \\
    \stage{Retrieve} & 0.00 & \textbf{1.00} & 1.00 & 0.00 & \textbf{1.00} & n.i. & n.i. \\
    \stage{Rank}     & 0.00 & 0.00 & \textbf{1.00} & 0.00 & 0.00 & n.i. & n.i. \\
    \stage{Explain}  & n.i. & n.i. & n.i. & 0.00 & 0.00 & 0.00 & n.i. \\
    \stage{Memory}   & n.i. & n.i. & n.i. & 0.00 & 0.00 & 0.00 & 0.00 \\
    \bottomrule
  \end{tabular}
\end{table}

The not-identified cells are a result rather than a gap in the table. A \stage{Rank} fault leaves
no measurable downstream footprint in a single turn: \stage{Explain}'s coordinate is a Jaccard
over reason tags and \stage{Memory}'s over stored facts, and neither reads the ranked list, so a
reordering is invisible to both. The same holds a fortiori for \stage{Explain} and \stage{Memory}
faults, which have one and zero downstream stages respectively. \textbf{In a single-turn
pipeline, repair removability is identified only for \stage{Elicit} and \stage{Retrieve}
faults}, and for the remaining three it is reported as not identified rather than filled
in---the same discipline Step 4 applies to harm when no defensible $\Omega_s$ exists. This
narrows Step 5 and is a further argument, alongside Section~\ref{sec:controlled-extensions}, for
treating \stage{Elicit}, \stage{Retrieve} and \stage{Rank} as the core.

\paragraph{Structural blindness separated from low power}
The masking result of Section~\ref{sec:controlled-step2} pools two claims of very different
strength; Section~\ref{sec:controlled-extensions} above separates them. Splitting the pairs by
whether the endpoint \emph{could} move: among the $80$ pairs planted at \stage{Elicit},
\stage{Retrieve} or \stage{Rank}, $5.0\%$ are process-sensitive yet endpoint-equivalent---a real
but modest power gap. The remaining $40$ pairs, planted at \stage{Explain} or \stage{Memory}, are
all process-sensitive and all endpoint-equivalent, with a maximum endpoint distance of
\emph{exactly} $0.0$. That is not a masking rate at all but the stronger structural claim, and it
is verified here rather than argued: an output-only audit is blind to \stage{Explain} and
\stage{Memory} faults by construction, not underpowered against them.

A follow-up turn removes the blindness. Under an explicit carry-over model---the agent's stored
facts and the reason tags it gave are echoed into the next turn's revealed preferences, since the
user continues from what the agent said and stored---every planted stage moves the turn-two
endpoint, including \stage{Explain} ($0.408$) and \stage{Memory} ($0.163$), while the controls
remain at exactly $0.000$. The blindness is a property of the single-turn horizon, not of the
stages, which is what makes the cross-turn treatment of \stage{Memory} load-bearing for any future
promotion of these two stages out of Section~\ref{sec:controlled-extensions} and into the core.

\begin{figure}[t]
  \centering
  \includegraphics[width=\linewidth]{fig_scaff_robustness.pdf}
  \caption{Three instrument properties that the prototype run left unestablished. (a) Candidate
    coordinates for \stage{Rank}, by their ability to separate \stage{Rank}-planted scenarios
    from all others when the crossover's noise is not held fixed; the dotted line is chance.
    (b) Localization as the decision margin varies, with the hand-set $\tau = 0.10$ marked: it
    sits above the effect scale of the recommended coordinate and costs $16.7$ points of accuracy
    on every seed. (c) Repairing the responsible stage against repairing a faithful one, for a
    fault planted at \stage{Elicit}, under the original scoring and under scoring restricted to
    stages strictly downstream of the repair.}
  \Description{Three panels. The left is a horizontal bar chart of detection AUC for five rank
    coordinates, with NDCG delta highest at about 0.80 and the top-one gap lowest at about 0.57,
    just above a dotted chance line at 0.5. The middle is a line chart of accuracy, recall and
    false-positive rate against the margin tau, with accuracy flat at about 0.83 across the
    region containing the hand-set margin and rising to 1.0 only as tau approaches zero. The
    right is a grouped bar chart with two groups; in the first group both bars are near 1.0, and
    in the second the true-stage bar remains at 1.0 while the faithful-stage bar drops to 0.2.}
  \label{fig:scaff-robustness}
\end{figure}

\subsection{What This Substrate Does and Does Not Establish}
\label{sec:controlled-limits}

It establishes that the crossover of Equation~\eqref{eq:crossovercell} recovers a known causal
origin, that natural divergence overattributes in a way the decomposition corrects, that direct
sensitivity is frequently invisible at the endpoint, and that repairing the implicated stage
removes substantially more downstream divergence than repairing an exonerated one. It does not
estimate any real-world disparity, its user oracle is exact by construction and therefore easier
than natural dialogue, its planted faults are stronger than the effects a deployed audit would
face, and it does not exercise the interactive multi-turn setting beyond a single follow-up turn.
The two evaluation settings the research questions call for---a controlled agentic
recommendation environment with verifiable constraints and repeated stochastic runs, and a
smaller external validation on a realistic interactive environment---are what these results are a
precondition for, not a substitute for.
